On veut neutraliser \qty{\v{Va}}{mL} d'une solution d'acide chlorhydrique (HCl) de concentration \qty{\v{Ca}}{mol/L} avec une solution d'hydroxyde de sodium (NaOH) de concentration \qty{\v{Cb}}{mol/L}. La réaction est :
\[ \mathrm{HCl + NaOH \longrightarrow NaCl + H_2O} \]

\begin{enumerate}[label=\alph*)]
    \item Combien de moles de HCl la solution d'acide contient-elle ?
    \item Quel volume de solution de NaOH faut-il ajouter pour neutraliser complètement l'acide ?
    \item Quelle masse de chlorure de sodium (NaCl, \qty{58.44}{g/mol}) est formée ?
\end{enumerate}
